\emph{(i)} At leading order in the patch size the metric at $p$ is flat and the approximate boost Killing field of a local Rindler horizon with edge $\Sigma$ is the flat boost that fixes $\Sigma$ (as in hypothesis (H2) of Theorem \claimref{synthesis.einstein-equations-from-c-loc-and-clausius}). In inertial coordinates $(t, x, y, z)$ at $p$ with $\Sigma = \{t = x = 0\}$ take $\ell = \partial_{t} + \partial_{x}$, $n = \partial_{t} - \partial_{x}$ and $K = x\,\partial_{t} + t\,\partial_{x}$. On the null hypersurface through $\Sigma$ generated by $\ell$, at the point $\lambda\ell$ with affine parameter $\lambda$, $K = \lambda\ell$; on the one generated by $n$, at $\lambda n$, $K = (-\lambda)\partial_{t} + \lambda\partial_{x} = -\lambda n$. On the past halves ($\lambda < 0$) $K$ is therefore past pointing on the $\ell$ sheet and future pointing on the $n$ sheet, and the future-directed boosts are $\chi_{\ell} = -K = -\lambda\ell$ and $\chi_{n} = K = -\lambda n$: the same field up to sign, each of the form $\chi = -\kappa\lambda k$ with $\kappa = 1$. On the future halves the orientations are reversed. The patches of a Killing horizon are local Rindler horizons (Theorem \claimref{synthesis.c-loc-is-a-second-placement-of-the-count}(ii)), so the same holds at a bifurcation surface with $K$ replaced by the Killing field. The check verifies the orientations and $|\chi| = |\lambda|$ on both halves of both sheets at three rational parameters.

\emph{(ii)} Fix a window $\mathfrak{A}_{F}$ of the shared algebra; it is a matrix factor, invariant under both flows, so the restriction of the state is KMS at the same period for the restricted flows (Proposition \claimref{hawking.kms-restricts-invariant-subalgebras}). Write $\alpha^{\ell}_{s} = \operatorname{Ad}(e^{JsH})$ on $\mathfrak{A}_{F}$; by (N), $\alpha^{n}_{s} = \operatorname{Ad}(e^{-JsH})$. By KMS uniqueness at a factor (Theorem \claimref{singularity.kms-uniqueness-factor}, with Lemma \claimref{hawking.gibbs-states-are-kms-conversely-finite}) the density is $e^{-2\pi H}/\operatorname{tr}e^{-2\pi H}$ and also $e^{2\pi H}/\operatorname{tr}e^{2\pi H}$. So $e^{-4\pi H}$ is a positive multiple of the identity; $H$ being self-adjoint, its eigenvalues are all equal, $H$ is a scalar, $\alpha^{\ell}$ is trivial on $\mathfrak{A}_{F}$ and the density is the identity divided by its trace, the normalised trace. Equivalently, by Theorem \claimref{singularity.rate-rigidity-temperature-not-algebra-intrinsic}(2), a flow-reversing identification carries the $2\pi$-KMS state to the $(-2\pi)$-KMS state, and on a factor these coincide only when the flow is trivial. Since every window is trivial for the flow and carries the trace, so does the inductive limit (the flow is trivial on a dense subalgebra, hence everywhere by continuity, and the state agrees with the trace on a dense subalgebra). For the phase-acting boost of the Definition, $H = \sum_{j \in F} w_{j}N_{j}$ on a window, scalar only if every weight $w_{j}$ vanishes. The check solves the linear KMS conditions $E^{-1}\varrho\,A\,E = A\varrho$ (all matrix units $A$, $E = e^{-\beta H}$ with a rational Boltzmann symbol) on a two-mode window with weights $1, 2$: one flow alone has a one-dimensional solution space, the Gibbs density; the flow and its inverse together have only $\varrho = 0$; a scalar $H$ admits the trace.

\emph{(iii)} On a window $\Mat_{n}(\C)$ every derivation is inner, and the derivations of a flow of $*$-automorphisms are $J[h, \cdot\,]$ with $h$ self-adjoint, determined up to a scalar; modulo scalars they form the Lie algebra $\mathfrak{su}(n)$ (identified with $J$ times the traceless self-adjoint matrices), on which $\operatorname{ad}$ of every element has purely imaginary eigenvalues, in $J\R$: $[Jh, E_{ab}] = J(\lambda_{a} - \lambda_{b})E_{ab}$ in an eigenbasis of $h$. Let $D$ be a Lie algebra homomorphism from the Lorentz algebra into it. In $\Cl(1,3)$ with the commutator bracket, the boost $B = e_{0}e_{1}$ has $[B, X_{\pm}] = \pm 2X_{\pm}$ for $X_{\pm} = e_{0}e_{2} \mp e_{1}e_{2}$ (computed in the check). Then $[D(B), D(X_{+})] = 2D(X_{+})$, so $D(X_{+})$ is zero or an eigenvector of $\operatorname{ad}_{D(B)}$ with the real nonzero eigenvalue $2$, which a purely imaginary spectrum excludes; hence $D(X_{+}) = 0$. The kernel of $D$ is an ideal containing $X_{+}$, and the ideal generated by $X_{+}$ is the whole six-dimensional bivector space (computed exactly in the check; this is the simplicity of the real Lie algebra $\mathfrak{so}(1,3)$, the realification of $\mathfrak{sl}(2,\C)$). So $D = 0$: a Lorentz-covariant action on finite windows is trivial whatever the state, and no condition on states, (R$_{\mathrm{loc}}$) included, is needed for this conclusion.

\emph{(iv)} By Theorem \claimref{hawking.modular-structure-completed-pair}(1) the completed pair is $v = \tfrac12(e_{0} + e_{1})$, $v^{*} = \tfrac12(e_{0} - e_{1})$ with $\sigma = e_{0}e_{1}$, and relabelling the pair reverses the orientation; $e_{0} \pm e_{1}$ are null and are the two generators $\ell$, $n$ of the sheets of the edge orthogonal to the plane of $\sigma$. A mode of weight $w$ has, for a flow $\operatorname{Ad}(e^{Jsw(N - 1/2)})$, the $2\pi$-KMS occupation $e^{-2\pi w}/(1 + e^{-2\pi w}) = 1/(1 + e^{2\pi w})$, and for the inverse flow $1/(1 + e^{-2\pi w})$; these agree only at $w = 0$. This is (ii) on the single-mode algebra $\Mat_{2}(\C)$.

\emph{(v)} In $\Cl(1,3)$ with $Q(e_{0}) = 1$, $Q(e_{i}) = -1$, the map $v \mapsto \tfrac12(e_{0}e_{1}v - v\,e_{0}e_{1})$ sends $te_{0} + xe_{1}$ to $-(xe_{0} + te_{1})$, so $B = e_{0}e_{1}$ is the bivector of $-K$, the future boost of the past half of the $\ell$ sheet, and $\tfrac12(B\ell - \ell B) = -\ell$. The edges orthogonal to $\ell$ are permuted transitively by the null rotations fixing $\ell$, the rotors $\Lambda_{a} = 1 + \tfrac{a}{2}\ell e_{j}$ ($j = 2, 3$; $(\ell e_{j})^{2} = 0$, $\Lambda_{a}\widetilde{\Lambda}_{a} = 1$, $\Lambda_{a}\ell\widetilde{\Lambda}_{a} = \ell$), which preserve the time orientation; the future boost of the horizon generated by $\ell$ with edge $\Lambda_{a}\Sigma$ is $B' = \Lambda_{a}B\widetilde{\Lambda}_{a}$, and $\tfrac12(B'\ell - \ell B') = \Lambda_{a}\tfrac12(B\ell - \ell B)\widetilde{\Lambda}_{a} = -\ell$. So $c_{\ell} = -1$ for every such boost. The stabiliser algebra of $[\ell]$ is spanned by $B$, $\ell e_{2}$, $\ell e_{3}$ and $e_{2}e_{3}$, with weights $-1, 0, 0, 0$ on $\ell$; $c_{\ell}$ is linear, and it vanishes on brackets because $\tfrac12\operatorname{ad}_{[X,Y]} = [\tfrac12\operatorname{ad}_{X}, \tfrac12\operatorname{ad}_{Y}]$ and both act on the line of $\ell$ by scalars; so it is a character. The check verifies, for three null rotations, that $B' \neq B$, $B'^{2} = 1$, that $B'$ moves $e_{j}$ (a different edge) and that its weight on $\ell$ is $-1$, the weights of the stabiliser basis, and that the brackets of that basis have weight zero; and that the opposite sheet's boost $-B$ gives $n$ the weight $-1$ and $\ell$ the weight $+1$. With the action $\alpha^{B}_{s} = \operatorname{Ad}(e^{-Jsc_{\ell}(B)h_{\ell}})$ on $\mathfrak{A}_{p,[\ell]}$ every horizon generated by $\ell$ acts on its modes at $p$ by the one flow $\operatorname{Ad}(e^{Jsh_{\ell}})$, whose $2\pi$-KMS state is $e^{-2\pi h_{\ell}}/\operatorname{tr}$; a horizon generated by $\ell$ contains no other direction's label at $p$, the algebras of different directions commute, and the product state restricts to each of these, so (R$_{\mathrm{loc}}$) holds for every local horizon through $p$ at every weight. At the other points $q$ of an edge the same construction applies with the null normal at $q$.

\emph{Consequences.} (a) Take as a hypothesis the reading of the stored Definition in which the hole's modes sit at the fixed locus of the Killing flow, the bifurcation surface, where Premise (R) asks each single-mode marginal to be KMS for the Killing flow with frequency $\omega_{j} > 0$ at $2\pi/\kappa$, rapidity period $2\pi$. A point of that surface is the edge of the two sheets of the bifurcate horizon, local Rindler horizons by (i) and the cited theorem, with future boosts $\xi$ and $-\xi$ on the past halves. Under that reading and the identification of (C$_{\mathrm{loc}}$) (one mode set) the hole's modes at the point are modes of both patches; each single-mode algebra is invariant, so (ii) applies to it with $H = \omega_{j}(N_{j} - \tfrac12)$ and gives $\omega_{j} = 0$, against the Definition. (b) In the conventions of Proposition \claimref{synthesis.de-sitter-static-boost-frozen-point}, the static Killing field of an observer is generated by an element of the stabiliser of the frozen point; in the embedding $-X_{0}^{2} + X_{1}^{2} + X_{2}^{2} + X_{3}^{2} + X_{4}^{2} = 1$ (unit radius) it is $\xi = X_{4}\partial_{0} + X_{0}\partial_{4}$, future pointing on the worldline $X_{4} > 0$ and past pointing on the antipodal worldline $X_{4} < 0$, and both observers' flows fix the sphere $X_{0} = X_{4} = 0$; the check verifies this at three rational points of each worldline and two of the sphere. If (R$_{\Lambda}$) is imposed for both observers on modes of that sphere, (ii) applies as in (a). (c) is (ii) and (v) together; (iii) shows that a covariant lift cannot replace the sheet labels. The per-direction windows of (v) are point spectrum, so by the stored no-go for half-sided translations on point spectrum the positive-null-energy premise of the result filed alongside (synthesis.positive-null-energy-gives-r-premises) cannot be posed on them; its continuum carrier is the other resolution.

\emph{Remarks.} The hypothesis that the patch's modes include modes at its edge is a reading of the stored Definition, not stored text; it is how the store places modes for a phase-acting boost: such a boost maps each mode to itself, and in the store the modes it acts on sit at its fixed locus ((R) at the bifurcation surface, the completion's pair over the exceptional set). On the continuum mode space with the boost dilating the affine parameter of the generators, modes away from the edge belong to one generator, so the labels of (v) are automatic there; in quantum field theory the algebras of the two sheets are those of the right and left wedges, which are each other's commutants (Bisognano and Wichmann; from memory). The null directions at a point form a two-sphere, which after complexification is the projective line of spinors $\pi$ with $\ell \propto \pi\bar\pi$, the incidence line of the point in the chart of Theorem \claimref{correspondence.incidence-relation-chart}; reading the labels of (v) as points of that line is a remark and is not used.
